\begin{abstract}
  A well-designed layout plays a crucial role in visual media, serving to both capture the audience's attention and facilitate information transfer. While this principle applies across various media, manga has a unique layout that reflects its narrative structure and is now a widely applied medium in fields beyond entertainment, such as education and advertising. In manga production, layouts are crafted based on a narrative script; consequently, the layout is a key component of visual direction, where the arrangement of panels and characters is constructed to effectively convey the narrative and enhance dramatic effect. Although layout generation has seen progress in domains such as Web UIs and posters that prioritize visual clarity, there has been little exploration into script-based layout generation, largely due to the lack of data. In this study, we created script data for approximately 20,000 pages from the Manga109 dataset and built the Manga109Script dataset linking scripts to manga layouts. Leveraging this dataset, we developed manga layout generation models that directly take scripts as input, employing various LLMs as base models. Our comprehensive evaluation shows that incorporating script information demonstrably improves layout generation quality compared to existing methods. These results establish an effective initial baseline for future research in script-based layout generation, which we term Script2Layout. As such, Script2Layout stands as a key enabling technology for applications designed to translate narratives into visual media, as exemplified by manga.
\end{abstract}
